\documentclass{article}

% Packages
\usepackage{geometry}
\usepackage{graphicx}
\usepackage{lipsum}
\usepackage{listings}
\usepackage[portuguese]{babel}

% Page setup
\geometry{a4paper, margin=2cm}

\begin{document}

% cover
\thispagestyle{empty} % Remove page number from first page

\begin{figure}[t]
    \includegraphics[width=3cm]{images/logo-puc-minas.png}
    \hspace{0.02\textwidth}
    \vline%
    \hspace{0.04\textwidth}
    \includegraphics[width=3cm]{images/logo-icei.jpeg}
\end{figure}

\hrulefill%
\vspace{\baselineskip}

\Large\noindent
\textbf{Pontifícia Universidade Católica de Minas Gerais} \\
\textbf{Instituto de Ciências Exatas e Informática} \\
\textbf{Departamento de Engenharia de Computação}

\begin{center}
    \vfill
    \Huge\textbf{Relatório: Trabalho Prático 2} \\
    \vspace{0.5cm}
    \Large\textbf{Registradores em VHDL} \\
    \vspace{1cm}
    \large \textbf{Professor(es)}: Antônio Hamilton Magalhães\\
    \vspace{0.5cm}
    \large \textbf{Aluno(s)}: Bruno Luiz Dias Alves de Castro \\
    \large \hspace{0.75cm} Rafael Ramos de Andrade \\
    \vfill
    \large Belo Horizonte \\ Campus Coração Eucarístico \\
    \vspace{\baselineskip}
    \large \today
\end{center}

% table of contents
\newpage
\thispagestyle{empty}
\tableofcontents

% body
\newpage
\large % document text size

\section{Introdução}

Durante as aulas da disciplina de Sistemas Reconfiguráveis, fomos introduzidos à linguagem VHDL. VHDL (\textbf{V}HSIC \textbf{H}ardware \textbf{D}escription \textbf{L}anguage) é uma linguagem de descrição de hardware. Com ela, podemos montar circuitos lógicos de maneira totalmente textual, o que garante à linguagem uma grande vantagem ante à soluções visuais.

\subsection{TP1}

Como primeiro trabalho prático (TP1), são propostas as montagens de dois circuitos utilizando a linguagem VHDL para a placa de desenvolvimento Altera:

\begin{itemize}
    \item Um ADDR\_MUX (Multiplexador de Endereçamento) (Secção 2)
    \item Uma ULA (Unidade Lógica Aritimética) (Secção 3)
\end{itemize}

Ambos os circuitos devem ser desenvolvidos utilizando programação concorrente, ou seja, sem trechos sequênciais no código-fonte.

\subsubsection{Objetivos}

Entre os objetivos que temos com o desenvolvimento deste trabalho prático podemos listar:

\begin{itemize}
    \item Aprender conceitos básicos da linguagem VHDL;
    \item Implementar utilizando programação concorrente os dois circuitos propostos;
    \item Compilar os circuitos e testar os resultados na placa de desenvolvimento Altera;
\end{itemize}

\section{ADDR\_MUX}
\lipsum[1]

\section{ULA (Unidade Lógica Aritimética)}

A ULA (\textbf{U}nidade \textbf{L}ógica \textbf{A}ritimétrica) é um dos componentes mais básicos de um processador. Como o nome já indica, a ULA é responsável por todas as operações lógicas (como OR, AND e Shift) e aritiméticas (como soma e subtração) realizadas em nosso circuito.

De maneira simplificada, a ULA receberá um comando, composto por seletores de operação e bits, e operandos. Na saída temos o resultado da operação desejada.

Como uma ULA opera de maneira concorrente, todas as operações implementadas são ``executadas ao mesmo tempo''. Um Multiplexador é usado para selecionar a operação correta.

\subsection{ULA Proposta}

Neste trabalho prático, a ULA proposta deve possuir as seguintes funções:

\begin{center}
    \begin{tabular}{|c|c|c|c|c|c|}
        \hline
        Função & op\_code & Descrição & z\_out & c\_out & dc\_out\\
        \hline
        XOR & 0000 & XOR Lógico & 1, se res. = 0 & - & -\\
        \hline
        OR & 0001 & OR Lógico & 1, se res. = 0 & - & -\\
        \hline
        AND & 0010 & AND Lógico & 1, se res. = 0 & - & -\\
        \hline
        CLR & 0011 & Limpa & 1 & - & -\\ 
        \hline
        ADD & 0100 & Soma & 1, se res. = 0 & 1 se \textit{carry} & 0 se \textit{carry} no nibble\\
        \hline
        SUB & 0101 & Subtração & 1, se res. = 0 & 0 se \textit{borrow} & 0 se \textit{borrow} no nibble\\
        \hline
        INC & 0110 & Incremento & 1, se res. = 0 & - & - \\
        \hline
        DEC & 0111 & Decremento & 1, se res. = 0 & - & - \\
        \hline
        PASS\_A & 1000 & Passa 'A' & 1, se res. = 0 & - & - \\
        \hline
        PASS\_B & 1001 & Passa 'B' & 1, se res. = 0 & - & - \\
        \hline
        COM & 1010 & Complemento & 1, se res. = 0 & - & - \\
        \hline
        SWAP & 1011 & Permuta \textit{nibbles} & 1, se res. = 0 & - & - \\
        \hline
        BS & 1100 & bit\_sel = 1 & a\_in[bit\_sel] & - & - \\
        \hline
        BC & 1101 & bit\_sel = 0 & a\_in[bit\_sel] & - & - \\
        \hline
        RR & 1110 & Rotação p/ dir. & - &  a\_in[0] & - \\
        \hline
        RL & 1111 & Rotação p/ esq. & - &  a\_in[7] & - \\
        \hline
    \end{tabular}
\end{center}

O sinais de entrada e saída são os seguintes:

\begin{center}
\begin{tabular}{|c|c|c|c|}
    \hline
    Nome & Tamanho & Tipo & Descrição\\
    \hline
    a\_in & 8 bits & \textit{Input} & Entrada de dados A\\
    \hline
    b\_in & 8 bits & \textit{Input} & Entrada de dados B\\
    \hline
    c\_in & 1 bit & \textit{Input} & Entrada de \textit{carry}\\
    \hline
    op\_sel & 4 bits & \textit{Input} & Seletor de operação\\
    \hline
    bit\_sel & 3 bits & \textit{Input} & Seletor de bit\\
    \hline
    r\_out & 8 bits & \textit{Output} & Saída do resultado\\
    \hline
    c\_out & 1 bit & \textit{Output} & Saída de \textit{carry/borrow}\\
    \hline
    dc\_out & 1 bit & \textit{Output} & Saída de \textit{digit carry/borrow}\\
    \hline
    z\_out & 1 bit & \textit{Output} & Saída de zero\\
    \hline
\end{tabular}
\end{center}

\subsection{Implementação}

Nossa implementação, feita em VHDL, consiste na utilização de métodos e bibloiotecas já implementadas, bem como lógica implementada por nós. Abaixo explicamos como foram implentadas cada uma das funções proportas.

\subsubsection{XOR, OR, AND e COM}

Esses métodos forma implemntados usando as funções lógicas disponíveis nativamente na linguagem VHDL.

\begin{lstlisting}[language=VHDL]
aux <=
    a XOR b WHEN "0000",    -- XOR
    a OR b  WHEN "0001",    -- OR
    a AND b WHEN "0010",    -- AND

    [...]

    NOT a   WHEN "1010", 	-- COM
\end{lstlisting}

\subsubsection{ADD, SUB, INC e DEC}

As funções de ADD (Adição), SUB (Subtração), INC (Incremento) e DEC (Decremento) foram implementadas usando operações aritiméticas já inclusas na linguagem. Para as funções de ADD e SUB, os operandos A e B são utilizados, e já para as operações de INC e DEC, apenas o operando A e a constante 1.

\begin{lstlisting}[language=VHDL]
aux <=  
    a + b   WHEN "0100",    -- ADD
    a - b   WHEN "0101",    -- SUB
    a + 1   WHEN "0110",    -- INC
    a - 1   WHEN "0111",    -- DEC
\end{lstlisting}

\subsubsection{PASS\_A e PASS\_B}

Talvez as operações mais simples, ambas não utilizam nenhuma lógica. As entradas A e B são apenas ``copiadas'' para a saída.

\begin{lstlisting}[language=VHDL]
aux <=  
    a   WHEN "1000",        -- PASS_A
    b   WHEN "1001",        -- PASS_B
\end{lstlisting}

\subsubsection{SWAP}

A operação de SWAP é feita invertando os dois \textit{nibbles}

\begin{lstlisting}[language=VHDL]
aux <= '0' & a(3 DOWNTO 0) & a(7 DOWNTO 4) WHEN "1011", -- SWAP
\end{lstlisting}

A concatenacão com zero no início será explicada na implementação do c\_out.

\subsubsection{BS e BC}

As operações de BS (\textit{bit set}) e BC (\textit{bit clear}) implementam a seguinte lógica: para ``setarmos'' um bit para 1, podemos fazer uma operação OR da entrada com uma \textit{string} de zeros, mas com um único um na posição que desejamos que seja igual à 1. Isso garanterá que este bit será sempre 1, e od demais são copiados. Podemos obter essa \textit{string} efetuando um \textit{shift lógico} de 1 um número bit\_sel de casas.

A operação BC funciona de maneira semalhando, mas deve ser feita com um AND no lugar do OR, e a \textit{string} deve ser de 1s com um único zero na posição desejada. Para isso, fazemos mesma operação que anteriormente, mas invertemos o resultado.

\begin{lstlisting}[language=VHDL]
aux <=
    a OR
        STD_LOGIC_VECTOR(SHIFT_LEFT(TO_UNSIGNED(1, 8),
            TO_INTEGER(UNSIGNED(bit_sel))))
        WHEN "1100",    -- BS

    a AND
        NOT STD_LOGIC_VECTOR(SHIFT_LEFT(TO_UNSIGNED(1, 8),
            TO_INTEGER(UNSIGNED(bit_sel))))
        WHEN "1101",	-- BC
\end{lstlisting}

\subsubsection{RR e RL}

As operações de RR (\textit{Rotate Right}) e RL (\textit{Rotate Left}) foram implementadas selecionando os bits necessários, e concatenando à entrada c\_in.

\begin{lstlisting}[language=VHDL]
aux <=
    '0' & c_in & a(7 DOWNTO 1) WHEN "1110",     -- RR
    '0' & a(6 DOWNTO 0) & c_in WHEN "1111";     -- RL
\end{lstlisting}

\subsubsection{z\_out}

A saída z\_out é usada para indicar que o resultado é zero para a maioria das operações, com duas excessões: a de CLR (já que o resultado sempre é zero), e as operações BS e BC onde ele deve ser o bit da entrada A na posição bit\_sel.

Para isso, verificamos se a operação termina com a sequência "110" (utilizada por ambas as operações) e colocamos o bit correto na saída. Caso contrário, ela é 1 caso a saída seja zero.

\begin{lstlisting}[language=VHDL]
z_out <=
    a(TO_INTEGER(UNSIGNED(bit_sel)))
        WHEN op_sel(3 DOWNTO 1) = "110" -- BS and BC

    ELSE '1' WHEN aux(7 DOWNTO 0) = "00000000";
\end{lstlisting}

\subsubsection{c\_out}

Para idetificarmos um \textit{carry} (ou um \textit{borrow}), há várias rotas que podemos tomar. Poderiamos, por exemplo, implementar cirtuitos lógicos capazar de identificar que ele ocorreu. Mas, para simplificar a operação, podemos chegar no mesmo resultado simplesmente aumentando o número de bits dos nossos dois operandos em 1. Podemos então copiar o bit mais significativo para a a saída c\_out. Caso um \textit{carry} ou um \textit{borrow} ocorra, esse bit será um.

\begin{lstlisting}[language=VHDL]
    r_out <= aux(7 DOWNTO 0);
    c_out <= aux(8);
\end{lstlisting}

Para evitarmos que o $9^o$ bit seja copiado para a saída, usamos um buffer aux, e apenas os bits de 7 à 0 são copiados para a saída c\_out.

\subsubsection{dc\_out}

Para identificarmos um \textit{carry} ou \textit{borrow} em um \textit{nibble}, podemos faer algo similar ao que fizemos c\_out. Criamos um auxiliar com 5 bits, e verificamos o copiamos para a saída, após executada a operação necessária.

\begin{lstlisting}[language=VHDL]
aux_nibble <=
    ('0' & a(3 DOWNTO 0)) +
        ('0' + b(3 DOWNTO 0)) WHEN op_sel(0) = '0' 
	ELSE ('0' & a(3 DOWNTO 0)) - ('0' & b(3 DOWNTO 0));

dc_out <= aux_nibble(4);
\end{lstlisting}

\section{Conclusion}
\lipsum[4]

\end{document}